\section{AHSME 1988}
        \begin{enumerate}
        	\item (\textit{AHSME 1988}) $\sqrt{8}+\sqrt{18}=$
$\text{(A)}\ \sqrt{20} \qquad  \text{(B)}\ 2(\sqrt{2}+\sqrt{3}) \qquad  \text{(C)}\ 7 \qquad  \text{(D)}\ 5\sqrt{2}\qquad \text{(E)}\ 2\sqrt{13}$

 \textbf{Answer: }D



\textbf{Solution}

$\sqrt{8} = 2  \sqrt{2}$
$\sqrt{18} = 3  \sqrt{2}$

 So adding the two terms we get $2  \sqrt{2} + 3  \sqrt{2} = 5  \sqrt{2}$ , which corresponds to answer choice $\boxed{\textbf{(D)}}$ .

 


 


	\item (\textit{AHSME 1988}) Triangles $ABC$ and $XYZ$ are similar, with $A$ corresponding to $X$ and $B$ to $Y$. If $AB=3, BC=4$, and $XY=5$, then $YZ$ is: 


 $\text{(A)}\ 3\frac{3}{4} \qquad  \text{(B)}\ 6 \qquad  \text{(C)}\ 6\frac{1}{4} \qquad  \text{(D)}\ 6\frac{2}{3} \qquad  \text{(E)}\ 8$

 \textbf{Answer: }D



\textbf{Solution}

Since the triangles are similar, we know that the corresponding sides of the triangle are in ratio to each other. We have $\frac{\overline{AB}}{\overline{XY}}=\frac{\overline{BC}}{\overline{YZ}}$. Plugging in values we have:$\frac{3}{4}=\frac{5}{\overline{YZ}}$. Solving for $\overline{YZ}$, we have $\overline{YZ}=6\frac{2}{3}$. So, the answer is $\boxed{\text{D}}$.

 


	\item (\textit{AHSME 1988}) 
                    \begin{figure}[H]
                        \centering
                        \includegraphics[width=0.5\textwidth]{Problem/AHSME_I_1988_3_Problem_Figure_1.png}
                    \end{figure}
                    


 Four rectangular paper strips of length $10$ and width $1$ are put flat on a table and overlap perpendicularly as shown. How much area of the table is covered?


 $\text{(A)}\ 36 \qquad  \text{(B)}\ 40 \qquad  \text{(C)}\ 44 \qquad  \text{(D)}\ 98 \qquad  \text{(E)}\ 100$

 \textbf{Answer: }A



\textbf{Solution}

We first notice that the paper strips cover up part of the others. Since the width of the overlap is $1$ and the length of the overlap is $1$, the area of each of the strips with the overlap is $(10\cdot 1)-1=9$. Since there are 4 strips, $4\cdot 9=36 \implies \boxed{\text{A}}$.

 


	\item (\textit{AHSME 1988}) The slope of the line $\frac{x}{3} + \frac{y}{2} = 1$ is


 $\textbf{(A)}\ -\frac{3}{2}\qquad \textbf{(B)}\ -\frac{2}{3}\qquad \textbf{(C)}\ \frac{1}{3}\qquad \textbf{(D)}\ \frac{2}{3}\qquad \textbf{(E)}\ \frac{3}{2}$

 \textbf{Answer: }B



\textbf{Solution}

To find the slope, all we have to do is put the equation into slope-intercept form. We subtract $\frac{x}{3}$ from both sides and then multiple all terms by $2$. This yields $y=-\frac{2}{3}x+1$, so the slope is $-\frac{2}{3} \implies \boxed{\text{B}}$.

 


	\item (\textit{AHSME 1988}) If $b$ and $c$ are constants and $(x + 2)(x + b) = x^2 + cx + 6$, then $c$ is


 $\textbf{(A)}\ -5\qquad \textbf{(B)}\ -3\qquad \textbf{(C)}\ -1\qquad \textbf{(D)}\ 3\qquad \textbf{(E)}\ 5$

 \textbf{Answer: }E



\textbf{Solution}

We first start out by expanding the left side of the equation, $(x+2)(x+b)=x^{2}+bx+2x+2b=x^2+(2+b)x+2b=x^2+cx+6$. We know the constant terms have to be equal so we have $2b=6$, so $b=3$. Plugging $b=3$ back in yields $x^2+(2+3)x+6=x^2+cx+6$. Thus, $c=5 \implies \boxed{\text{E}}$.

 


	\item (\textit{AHSME 1988}) A figure is an equiangular parallelogram if and only if it is a


 $\textbf{(A)}\ \text{rectangle}\qquad \textbf{(B)}\ \text{regular polygon}\qquad \textbf{(C)}\ \text{rhombus}\qquad \textbf{(D)}\ \text{square}\qquad \textbf{(E)}\ \text{trapezoid}$

 \textbf{Answer: }A



\textbf{Solution}

The definition of an equiangular parallelogram is that all the angles are equal, and that pairs of sides are parallel. It may be a rectangle, because all the angles are equal and it is a parallelogram. It is not necessarily a regular polygon, because if the polygon is a pentagon, it is not a parallelogram. It is not necessarily a rhombus, because all the angles are not necessarily equal.  It may be a square, since it is a parallelogram and all the angles are equal. It is not necessarily a trapezoid, because the angles are not necessarily equal. We have that it could be a square or a rectangle. A square is a rectangle, but a rectangle is not necessarily a square. We want the all-encompassing answer so the answer is a rectange $\implies \boxed{\text{A}}$.

 


	\item (\textit{AHSME 1988}) Estimate the time it takes to send $60$ blocks of data over a communications channel if each block consists of $512$
"chunks" and the channel can transmit $120$ chunks per second.


 $\textbf{(A)}\ 0.04 \text{ seconds}\qquad \textbf{(B)}\ 0.4 \text{ seconds}\qquad \textbf{(C)}\ 4 \text{ seconds}\qquad \textbf{(D)}\ 4\text{ minutes}\qquad \textbf{(E)}\ 4\text{ hours}$

 \textbf{Answer: }D



\textbf{Solution 1}

We want to figure out the number of chunks in $60$ blocks, so we have $60\cdot 512 \approx 30000$. We divide this by $120$ to determine the number of seconds necessary to transmit. $30000/120 \approx 250$, which means that it takes approximately $4$ minutes to transmit. Thus, the answer is $\boxed{\text{D}}$.

 



\textbf{Solution 2}

This solution is if you are running out of time and just want to write down an answer. So, this is quite unreliable. You can logic it out. It doesn't make sense for the first three options to be the answer since that is way too quick. The last option is way too long. That just leaves $\boxed{\text{D}}$.

 


	\item (\textit{AHSME 1988}) If $\frac{b}{a} = 2$ and $\frac{c}{b} = 3$, what is the ratio of $a + b$ to $b + c$?


 $\textbf{(A)}\ \frac{1}{3}\qquad \textbf{(B)}\ \frac{3}{8}\qquad \textbf{(C)}\ \frac{3}{5}\qquad \textbf{(D)}\ \frac{2}{3}\qquad \textbf{(E)}\ \frac{3}{4}$

 \textbf{Answer: }B



\textbf{Solution 1}

Since we are finding ratios, it would be helpful to put everything in terms of one variable. Since $b$ is in both equations, that would be a place to start. We manipulate the equations yielding $\frac{b}{2}=a$ and $c=3b$. Since we are asked to find the ratio of $a+b$ to $b+c$, we need to find $\frac{a+b}{b+c}$. We found the $a$ and $c$ in terms of $b$ so that means we can plug them in. We have: $\frac{\frac{b}{2}+b}{b+3b}=\frac{\frac{3}{2}b}{4b}=\frac{3}{8}$. Thus the answer is $\frac{3}{8} \implies \boxed{\text{B}}$.

 



\textbf{Solution 2}

WLOG, let $b=4, a=2, c=12.$Thus, the answer is $\frac{4+2}{12+4}= \frac{3}{8}$



	\item (\textit{AHSME 1988}) 
                    \begin{figure}[H]
                        \centering
                        \includegraphics[width=0.5\textwidth]{Problem/AHSME_I_1988_9_Problem_Figure_1.png}
                    \end{figure}
                    


 An $8' \times 10'$ table sits in the corner of a square room, as in Figure $1$ below. 
The owners desire to move the table to the position shown in Figure $2$. 
The side of the room is $S$ feet. What is the smallest integer value of $S$ for which the table can be moved as desired without tilting it or taking it apart?


 $\textbf{(A)}\ 11\qquad \textbf{(B)}\ 12\qquad \textbf{(C)}\ 13\qquad \textbf{(D)}\ 14\qquad \textbf{(E)}\ 15$

 \textbf{Answer: }C



\textbf{Solution}

Small correction: The writer below has maximized the area of the rectangle (with sides parallel to the walls) that fits around the table, but there is a larger single dimension we can find in the table. The height or width is maximized when the diagonal of the table is horizontal or vertical. By the Pythagorean Theorem, this diagonal is $\sqrt{8^2+10^2} = \sqrt{164},$ which is between $\sqrt{144}$ and $\sqrt{169},$ so the answer is still $\textbf{(C)}.$ -hailstone

 We begin by thinking about the motion of the table. As it moves, the table will have it's maximum height and width when the rectangle's sides form $45$ degree angles relative to the sides of the square. Therefore, by the Pythagorean Theorem, we have that $S= 4\sqrt{2}+5\sqrt{2}$, with $4\sqrt{2}$ being the length of the leg formed by the side of the square with length $8$ and $5\sqrt{2}$ being the length of the leg formed by the side of the square with length $10$. Adding these up yields $9\sqrt{2}$.

 We have that $\sqrt{2}\approx 1.414\approx 1.4$. That means that $9\sqrt{2}\approx 12.6$, which rounds up to $\boxed{\textbf{C)} 13}$.

 


	\item (\textit{AHSME 1988}) In an experiment, a scientific constant $C$ is determined to be $2.43865$ with an error of at most $\pm 0.00312$. 
The experimenter wishes to announce a value for $C$ in which every digit is significant. 
That is, whatever $C$ is, the announced value must be the correct result when $C$ is rounded to that number of digits. 
The most accurate value the experimenter can announce for $C$ is


 $\textbf{(A)}\ 2\qquad \textbf{(B)}\ 2.4\qquad \textbf{(C)}\ 2.43\qquad \textbf{(D)}\ 2.44\qquad \textbf{(E)}\ 2.439$

 \textbf{Answer: }D



\textbf{Solution}

If added together, we have:
 \[2.43865+0.00312=2.44177\]This rounds to $2.44$. If they subtracted, we have:
 \[2.43865-0.00312=2.43553.\]This rounds to $2.44$. Therefore, we have the answer to be $\fbox{\textbf{(D)} 2.44}$.

 


	\item (\textit{AHSME 1988}) 
                    \begin{figure}[H]
                        \centering
                        \includegraphics[width=0.5\textwidth]{Problem/AHSME_I_1988_11_Problem_Figure_1.png}
                    \end{figure}
                    


 On each horizontal line in the figure below, the five large dots indicate the populations of cities $A, B, C, D$ and $E$ in the year indicated. 
Which city had the greatest percentage increase in population from $1970$ to $1980$?


 $\textbf{(A)}\ A\qquad \textbf{(B)}\ B\qquad \textbf{(C)}\ C\qquad \textbf{(D)}\ D\qquad \textbf{(E)}\ E$

 \textbf{Answer: }C



\textbf{Solution}

We simply have to divide the present population by the original, and whichever one is the largest have the greatest percent change. City $\text{A}$  is $\frac{50}{40}=1.25$. City $\text{B}$ is $\frac{70}{50}=1.4$. City $\text{C}$ is $\frac{100}{70} \approx 1.42$. City $\text{D}$ is $\frac{130}{100}=1.3$. City $\text{E}$ is $\frac{160}{120} \approx 1.33$. City $\text{C}$ has the greatest value, so the answer is $\boxed{\text{C}}$.

 


	\item (\textit{AHSME 1988}) Each integer $1$ through $9$ is written on a separate slip of paper and all nine slips are put into a hat. 
Jack picks one of these slips at random and puts it back. Then Jill picks a slip at random. 
Which digit is most likely to be the units digit of the sum of Jack's integer and Jill's integer?


 $\textbf{(A)}\ 0\qquad \textbf{(B)}\ 1\qquad \textbf{(C)}\ 8\qquad \textbf{(D)}\ 9\qquad \textbf{(E)}\ \text{each digit is equally likely}$

 \textbf{Answer: }A



\textbf{Solution}

We can draw a sample space diagram, and find that of the $9^2 = 81$ possibilities, $9$ of them give a sum of $0$, and each other sum mod $10$ (from $1$ to $9$) is given by $8$ of the possibilities (and indeed we can check that $9 + 8 \times 9 = 81$). Thus $0$ is the most likely, so the answer is $\boxed{\text{A}}$.

 



\textbf{Solution (alternative explanation)}

The largest sum that can be formed from two slips of these papers is 18, and if we list the composition of each of the numbers greater than 1 and less than 19, we can find that there are 9 combinations in total which make a sum of 10 ($1+9$, $2+8$, etc). Notice that the order in which either number appears doesn't matter since they all count toward the number of possibilities.

 


	\item (\textit{AHSME 1988}) If $\sin(x)= 3\cos(x)$ then what is $\sin(x) \cdot \cos(x)$?


 $\textbf{(A)}\ \frac{1}{6}\qquad \textbf{(B)}\ \frac{1}{5}\qquad \textbf{(C)}\ \frac{2}{9}\qquad \textbf{(D)}\ \frac{1}{4}\qquad \textbf{(E)}\ \frac{3}{10}$

 \textbf{Answer: }E



\textbf{Solution}

In the problem we are given that $\sin{(x)}=3\cos{(x)}$, and we want to find $\sin{(x)}\cos{(x)}$. We can divide both sides of the original equation by $\cos{(x)}$ to get 
 \[\frac{\sin{(x)}}{\cos{(x)}}=\tan{(x)}=3.\]We can now use right triangle trigonometry to finish the problem.
                    \begin{figure}[H]
                        \centering
                        \includegraphics[width=0.5\textwidth]{Solution/AHSME_I_1988_13_Solution_Figure_1.png}
                    \end{figure}
                    

 (Note that this assumes that $x$ is acute. If $x$ is obtuse, then $\sin{(x)}$ is positive and $\cos{(x)}$ is negative, so the equation cannot be satisfied. If $x$ is reflex, then both $\sin{(x)}$ and $\cos{(x)}$ are negative, so the equation is satisfied, but when we find $\sin{(x)}\cos{(x)}$, the two negatives will cancel out and give the same (positive) answer as in the acute case.)

 Since the problem asks us to find $\sin{(x)}\cos{(x)}$.
 \[\sin{(x)}\cos{(x)}=\left(\frac{3}{\sqrt{10}}\right)\left(\frac{1}{\sqrt{10}}\right)=\frac{3}{10}.\]So $\boxed{\textbf{(E)}\ \frac{3}{10}}$ is our answer.

 



\textbf{Solution 2 (Pure Algebra)}

Squaring both sides gives ${\sin}^2 x = 9{\cos}^2 x$. We can take the Pythagorean identity, ${\sin}^2 x + {\cos}^2 x = 1$ and substitute the 1st equation in, giving $10{\cos}^2 x = 1$. So ${\cos}^2 x = \frac{1}{10}$, and ${\sin}^2 x = \frac{9}{10}$. 

 Multiplying the 2 together gives ${\sin}^2 x {\cos}^2 x = \frac{9}{100}$, and then taking the square root gives $\mp \frac{3}{10}$. However, $-\frac{3}{10}$ implies one of $\sin x$ and $\cos x$ is negative while the other is positive, but $\sin x = 3 \cos x$ means they have the same sign, which contradicts the first statement. This means $\boxed{\textbf{(E)}\ \frac{3}{10}}$ is the only answer.

 -ThisUsernameIsTaken

 


	\item (\textit{AHSME 1988}) For any real number a and positive integer k, define


 ${a \choose k} = \frac{a(a-1)(a-2)\cdots(a-(k-1))}{k(k-1)(k-2)\cdots(2)(1)}$


 What is


 ${-\frac{1}{2} \choose 100} \div {\frac{1}{2} \choose 100}$?


 $\textbf{(A)}\ -199\qquad \textbf{(B)}\ -197\qquad \textbf{(C)}\ -1\qquad \textbf{(D)}\ 197\qquad \textbf{(E)}\ 199$

 \textbf{Answer: }A



\textbf{Solution}

We expand both the numerator and the denominator.

 
 \begin{align*} \binom{-\frac{1}{2}}{100}\div\binom{\frac{1}{2}}{100} &= \frac{   \dfrac{     (-\frac{1}{2})     (-\frac{1}{2} - 1)     (-\frac{1}{2} - 2)     \cdots     (-\frac{1}{2} - (100 - 1))   }{\cancel{(100)(99)\cdots(1)}} }{   \dfrac{     (\frac{1}{2})     (\frac{1}{2} - 1)     (\frac{1}{2} - 2)     \cdots     (\frac{1}{2} - (100 - 1))   }{\cancel{(100)(99)\cdots(1)}} } \\ &= \frac{     (-\frac{1}{2})     (-\frac{1}{2} - 1)     (-\frac{1}{2} - 2)     \cdots     (-\frac{1}{2} - 99) }{     (\frac{1}{2})     (\frac{1}{2} - 1)     (\frac{1}{2} - 2)     \cdots     (\frac{1}{2} - 99) } \end{align*}


 Now, note that $-\frac{1}{2}-1=\frac{1}{2}-2$, $-\frac{1}{2}-2=\frac{1}{2}-3$, etc.; in essence, $-\frac{1}{2}-n=\frac{1}{2}-(n+1)$. We can then simplify the numerator and cancel like terms.

 
 \begin{align*} \frac{     (-\frac{1}{2})     (-\frac{1}{2} - 1)     (-\frac{1}{2} - 2)     \cdots     (-\frac{1}{2} - 99) }{     (\frac{1}{2})     (\frac{1}{2} - 1)     (\frac{1}{2} - 2)     \cdots     (\frac{1}{2} - 99) } &= \frac{     \cancel{(\frac{1}{2} - 1)}     \cancel{(\frac{1}{2} - 2)}     \cancel{(\frac{1}{2} - 3)}     \cdots     (\frac{1}{2} - 100) }{     (\frac{1}{2})     \cancel{(\frac{1}{2} - 1)}     \cancel{(\frac{1}{2} - 2)}     \cdots     \cancel{(\frac{1}{2} - 99)} } \\ &= \frac{\frac{1}{2}-100}{\frac{1}{2}} \\ &= \frac{-\frac{199}{2}}{\frac{1}{2}} \\ &= \boxed{\textbf{(A) } -199.} \end{align*}




	\item (\textit{AHSME 1988}) If $a$ and $b$ are integers such that $x^2 - x - 1$ is a factor of $ax^3 + bx^2 + 1$, then $b$ is


 $\textbf{(A)}\ -2\qquad \textbf{(B)}\ -1\qquad \textbf{(C)}\ 0\qquad \textbf{(D)}\ 1\qquad \textbf{(E)}\ 2$

 \textbf{Answer: }A



\textbf{Solution}

Using polynomial division, we find that the remainder is $(2a+b)x+(a+b+1)$, so for the condition to hold, we need this remainder to be $0$. This gives $2a+b=0$ and $a+b+1=0$, so $b=-2a$ and $a-2a+1=0 \implies a=1 \implies b=-2$, which is $\boxed{\text{A}}.$



	\item (\textit{AHSME 1988}) 
                    \begin{figure}[H]
                        \centering
                        \includegraphics[width=0.5\textwidth]{Problem/AHSME_I_1988_16_Problem_Figure_1.png}
                    \end{figure}
                    
$ABC$ and $A'B'C'$ are equilateral triangles with parallel sides and the same center, 
as in the figure. The distance between side $BC$ and side $B'C'$ is $\frac{1}{6}$ the altitude of $\triangle ABC$. 
The ratio of the area of $\triangle A'B'C'$ to the area of $\triangle ABC$ is


 $\textbf{(A)}\ \frac{1}{36}\qquad \textbf{(B)}\ \frac{1}{6}\qquad \textbf{(C)}\ \frac{1}{4}\qquad \textbf{(D)}\ \frac{\sqrt{3}}{4}\qquad \textbf{(E)}\ \frac{9+8\sqrt{3}}{36}$

 \textbf{Answer: }C



\textbf{Solution}

Let $\triangle ABC$ have side length $s$ and $\triangle A'B'C'$ have side length $t$. Thus the altitude of $\triangle ABC$ is $\frac{s\sqrt{3}}{2}$. Now observe that this altitude is made up of three parts: the distance from $BC$ to $B'C'$, plus the altitude of $\triangle A'B'C'$, plus a top part which is equal to the length of the diagonal line from the bottom-left corner of $\triangle ABC$ to the bottom left corner of $\triangle A'B'C'$ (as an isosceles trapezium is formed with parallel sides $AB$ and $A'B'$, and legs $AA'$ and $BB'$). We drop a perpendicular from $B'$ to $BC$, which meets $BC$ at $D$. $\triangle BDB'$ has angles $30^{\circ}$, $60^{\circ}$, and $90^{\circ}$, and the vertical side is that distance from $BC$ to $B'C'$, which is given as $\frac{1}{6} \times \frac{s\sqrt{3}}{2} = \frac{s\sqrt{3}}{12}$, so that by the length relationships in a 30-60-90 triangle,, the length of the diagonal line is $\frac{s\sqrt{3}}{12} \times 2 = \frac{s\sqrt{3}}{6}.$ Thus using the "altitude in three parts" idea, we get $\frac{s\sqrt{3}}{2} = \frac{s\sqrt{3}}{12} + \frac{t\sqrt{3}}{2} + \frac{s\sqrt{3}}{6} \implies \frac{s\sqrt{3}}{4} = \frac{t\sqrt{3}}{2} \implies t = \frac{1}{2}s.$ Thus the sides of $\triangle A'B'C'$ are half as long as $\triangle ABC$, so the area ratio is $(\frac{1}{2}) ^ {2} = \frac{1}{4}$, which is $\boxed{\text{C}}$.

 


	\item (\textit{AHSME 1988}) If $|x| + x + y = 10$ and $x + |y| - y = 12$, find $x + y$
$\textbf{(A)}\ -2\qquad \textbf{(B)}\ 2\qquad \textbf{(C)}\ \frac{18}{5}\qquad \textbf{(D)}\ \frac{22}{3}\qquad \textbf{(E)}\ 22$

 \textbf{Answer: }C



\textbf{Solution}

We proceed by casework:

 Case 1: x, y >= 0

 Thus we have 2x + y = 10, x = 12. This makes y < 0, a contradiction.

 Case 2: x >= 0 > y

 2x + y = 10, x - 2y = 12. By multiplying the first equation by two and adding, we have 4x + x = 2*10 + 12 thus x = 32/5. Since the question only specifies 1 solution, plugging in x = 32/5 into the first equation yields y = -14/5 thus x + y = 18/5 $C$.

 


	\item (\textit{AHSME 1988}) At the end of a professional bowling tournament, the top 5 bowlers have a playoff. 
First \#5 bowls \#4. The loser receives $5$th prize and the winner bowls \#3 in another game. 
The loser of this game receives $4$th prize and the winner bowls \#2. 
The loser of this game receives $3$rd prize and the winner bowls \#1. 
The winner of this game gets 1st prize and the loser gets 2nd prize. 
In how many orders can bowlers \#1 through \#5 receive the prizes?


 $\textbf{(A)}\ 10\qquad \textbf{(B)}\ 16\qquad \textbf{(C)}\ 24\qquad \textbf{(D)}\ 120\qquad \textbf{(E)}\ \text{none of these}$

 \textbf{Answer: }B



\textbf{Solution}

We have $2$ choices for who wins the first game, and that uniquely determines $5^{\text{th}}$ place. Then there are $2$ choices for a next game and that uniquely determines $4^{\text{th}}$ place, followed by $2$ choices for the next game that uniquely determines $3^{\text{rd}}$ place. Finally, there are $2$ choices for the last game, which uniquely determines both $1^{\text{st}}$ and $2^{\text{nd}}$ places, since the winner is $1^{\text{st}}$ and the loser is $2^{\text{nd}}$. Thus the number of possible orders is $2 \times 2 \times 2 \times 2 = 16$, which is $\boxed{\text{B}}$.

 See Volume I Chapter 25, p. 234-235

 


	\item (\textit{AHSME 1988}) Simplify


 $\frac{bx(a^2x^2 + 2a^2y^2 + b^2y^2) + ay(a^2x^2 + 2b^2x^2 + b^2y^2)}{bx + ay}$
$\textbf{(A)}\ a^2x^2 + b^2y^2\qquad \textbf{(B)}\ (ax + by)^2\qquad \textbf{(C)}\ (ax + by)(bx + ay)\qquad\\ \textbf{(D)}\ 2(a^2x^2+b^2y^2)\qquad \textbf{(E)}\ (bx+ay)^2$

 \textbf{Answer: }B



\textbf{Solution 1}

We can multiply each answer choice by $bx + ay$ and then compare with the numerator. This gives $\boxed{\text{B}}$.

 



\textbf{Solution 2}

Expanding everything in the brackets, we get$\frac{ba^2x^3 + 2ba^2xy^2 + b^3xy^2 + a^3x^2y + 2ab^2x^2y + ab^2y^3}{bx + ay}$. We can then group numbers up in pairs so they equal $n(bx + ay)$:

 $= \frac{ba^2x^3 + a^3x^2y + 2ab^2x^2y + 2ba^2xy^2 + b^3xy^2 + ab^2y^3}{bx+ay}$
$= \frac{bx + ay(a^2x^2) + bx + ay(2baxy) + bx + ay(b^2y^2)}{bx+ay}$
$= a^2x^2 + 2baxy + b^2y^2$
$= (ax + by)^2$

 We get $\boxed{\text{B}}$.

 -ThisUsernameIsTaken

 



\textbf{Solution 3}

If you were out of time and your algebra isn't that good, you could just plug in some values for the variables and see which answer choice works.

 



\textbf{Solution 4 (fastest)}

After regrouping, the numerator becomes $(bx+ay)(a^2x^2+b^2y^2)+2bxa^2y^2+2ayb^2x^2$. Factoring further, we get $(bx+ay)(a^2x^2+b^2y^2)+2bxay(bx+ay)$. After dividing, we get $a^2x^2+b^2y^2+2bxay$, which can be factored as $(ax+by)^2$, so the answer is $\boxed{\text{B}}$.

 -Pengu14

 


	\item (\textit{AHSME 1988}) In one of the adjoining figures a square of side $2$ is dissected into four pieces so that $E$ and $F$ are the midpoints
of opposite sides and $AG$ is perpendicular to $BF$. These four pieces can then be reassembled into a rectangle as shown 
in the second figure. The ratio of height to base, $XY / YZ$, in this rectangle is  


 
                    \begin{figure}[H]
                        \centering
                        \includegraphics[width=0.5\textwidth]{Problem/AHSME_I_1988_20_Problem_Figure_1.png}
                    \end{figure}
                    
$\textbf{(A)}\ 4\qquad \textbf{(B)}\ 1+2\sqrt{3}\qquad \textbf{(C)}\ 2\sqrt{5}\qquad \textbf{(D)}\ \frac{8+4\sqrt{3}}{3}\qquad \textbf{(E)}\ 5$

 \textbf{Answer: }E



\textbf{Solution}

Within $WXYZ$, the parallelogram piece has vertical side $BF = \sqrt{1^2 + 2^2} = \sqrt{5}$, and diagonal side $FD = 1$. Thus the triangle in the bottom-right hand corner (the one with horizontal side $YZ$) must have hypotenuse $1$, and the only such triangle in the original figure is $\triangle AFG$, so we deduce $YZ = AG = \frac{\frac{1}{2} \times 1 \times 2}{\frac{1}{2} \times \sqrt{5}} = \frac{2}{\sqrt{5}}.$ Now the rectangle must have the same area as the square, as the pieces were put together without gaps or overlaps, so its area is $2^2 = 4$, and thus the vertical side $WZ$ is $\frac{4}{\frac{2}{\sqrt{5}}} = 2\sqrt{5}$, so the required ratio is $\frac{2\sqrt{5}}{\frac{2}{\sqrt{5}}} = \sqrt{5} \times \sqrt{5} = 5$, which is $\boxed{\text{E}}$.

 


	\item (\textit{AHSME 1988}) The complex number $z$ satisfies $z + |z| = 2 + 8i$. What is $|z|^{2}$? Note: if $z = a + bi$, then $|z| = \sqrt{a^{2} + b^{2}}$.


 $\textbf{(A)}\ 68\qquad \textbf{(B)}\ 100\qquad \textbf{(C)}\ 169\qquad \textbf{(D)}\ 208\qquad \textbf{(E)}\ 289$

 \textbf{Answer: }E



\textbf{Solution}

Let the complex number $z$ equal $a+bi$. Then the preceding equation can be expressed as 
 \[a+bi+\sqrt{a^2+b^2} = 2+8i\] Because $a$ and $b$ must both be real numbers, we immediately have that $bi = 8i$, giving $b = 8$. Plugging this in back to our equation gives us $a+\sqrt{a^2+64} = 2$.  Rearranging this into $2-a = \sqrt{a^2+64}$, we can square each side of the equation resulting in 
 \[4-4a+a^2 = a^2+64\] Further simplification will yield  $60 = -4a$ meaning that $-15 = a$. Knowing both $a$ and $b$, we can plug them in into $a^2+b^2$. Our final answer is $\boxed{289}$.

 


	\item (\textit{AHSME 1988}) For how many integers $x$ does a triangle with side lengths $10, 24$ and $x$ have all its angles acute?


 $\textbf{(A)}\ 4\qquad \textbf{(B)}\ 5\qquad \textbf{(C)}\ 6\qquad \textbf{(D)}\ 7\qquad \textbf{(E)}\ \text{more than } 7$

 \textbf{Answer: }A



\textbf{Solution}

We first notice that the sides $10$ and $24$, can be part of $2$ different right triangles, one with sides $10,24,26$, and the other with a legsomewhere between $21$ and $22$. We now notice that if $x$ is less than or equal to $21$, one of the angles is obtuse, and that the same is the same for any value of $x$ above $26$. Thus the only integer values of $x$ that fit the conditions, are $x=22, 23, 24, \text{and }25.$So, the answer is $\boxed{\text{A}}$



	\item (\textit{AHSME 1988}) The six edges of a tetrahedron $ABCD$ measure $7, 13, 18, 27, 36$ and $41$ units. If the length of edge $AB$ is $41$, then the length of edge $CD$ is


 $\textbf{(A)}\ 7\qquad \textbf{(B)}\ 13\qquad \textbf{(C)}\ 18\qquad \textbf{(D)}\ 27\qquad \textbf{(E)}\ 36$

 \textbf{Answer: }B



\textbf{Solution}

By the triangle inequality in $\triangle ABC$, we find that $BC$ and $CA$ must sum to greater than $41$, so they must be (in some order) $7$ and $36$, $13$ and $36$, $18$ and $27$, $18$ and $36$, or $27$ and $36$. We try $7$ and $36$, and now by the triangle inequality in $\triangle ABD$, we must use the remaining numbers $13$, $18$, and $27$ to get a sum greater than $41$, so the only possibility is $18$ and $27$. This works as we can put $BC = 36$, $AC = 7$, $AD = 18$, $BD = 27$, $CD = 13$, so that $\triangle ADC$ and $\triangle BDC$ also satisfy the triangle inequality. Hence we have found a solution that works, and it can be verified that the other possibilities don't work, though as this is a multiple-choice competition, you probably wouldn't do that in order to save time. In any case, the answer is $CD = 13$, which is $\boxed{\text{B}}$.

 


 


	\item (\textit{AHSME 1988}) An isosceles trapezoid is circumscribed around a circle. The longer base of the trapezoid is $16$, 
and one of the base angles is $\arcsin(.8)$. Find the area of the trapezoid.


 $\textbf{(A)}\ 72\qquad \textbf{(B)}\ 75\qquad \textbf{(C)}\ 80\qquad \textbf{(D)}\ 90\qquad \textbf{(E)}\ \text{not uniquely determined}$

 \textbf{Answer: }C



\textbf{Solution}

Let the trapezium have diagonal legs of length $x$ and a shorter base of length $y$. Drop altitudes from the endpoints of the shorter base to the longer base to form two right-angled triangles, which are congruent since the trapezium is isosceles. Thus using the base angle of $\arcsin(0.8)$ gives the vertical side of these triangles as $0.8x$ and the horizontal side as $0.6x$. Now notice that the sides of the trapezium can be seen as being made up of tangents to the circle, and thus using the fact that "the tangents from a point to a circle are equal in length" gives $2y + 0.6x + 0.6x = 2x$. Also, using the given length of the longer base tells us that $y + 0.6x + 0.6x = 16$. Solving these equations simultaneously gives $x=10$ and $y=4$, so the height of the trapezium is $0.8 \times 10 = 8$. Thus the area is $\frac{1}{2}(4+16)(8) = 80$, which is $\boxed{\text{C}}$.

 


 


	\item (\textit{AHSME 1988}) $X, Y$ and $Z$ are pairwise disjoint sets of people. The average ages of people in the sets 
$X, Y, Z, X \cup Y, X \cup Z$ and $Y \cup Z$ are $37, 23, 41, 29, 39.5$ and $33$ respectively. 
Find the average age of the people in set $X \cup Y \cup Z$.


 $\textbf{(A)}\ 33\qquad \textbf{(B)}\ 33.5\qquad \textbf{(C)}\ 33.6\overline{6}\qquad \textbf{(D)}\ 33.83\overline{3}\qquad \textbf{(E)}\ 34$

 \textbf{Answer: }E



\textbf{Solution}

Let the variables $X$, $Y$, and $Z$ represent the sums of the ages of the people in sets $X$, $Y$, and $Z$ respectively. Let $x$, $y$, and $z$ represent the numbers of people who are in sets $X$, $Y$, and $Z$ respectively. Since the sets are disjoint, we know, for example, that the number of people in the set $X \cup Y$ is $x+y$ and the sum of their ages is $X+Y$, and similar results apply for the other unions of sets. Thus we have $X=37x$, $Y=23y$, $Z=41z$, $X+Y=29(x+y) \implies 37x + 23y = 29x + 29y \implies 8x = 6y \implies y = \frac{4}{3}x$, and $X+Z = 39.5(x+z) \implies 37x + 41z = 39.5x + 39.5z \implies 74x + 82z = 79x + 79z \implies 3z = 5x \implies z = \frac{5}{3}x$. Hence the answer is $\frac{X+Y+Z}{x+y+z} = \frac{37x + 23\times\frac{4}{3}x + 41\times\frac{5}{3}x}{x+\frac{4}{3}x+\frac{5}{3}x} = \frac{136x}{4x} = 34$, which is $\boxed{\text{E}}$. Notice that we did not need all three of the means of the unions of the sets - any two of them would have been sufficient to determine the answer.

 


 


	\item (\textit{AHSME 1988}) Suppose that $p$ and $q$ are positive numbers for which


 $\log_{9}(p) = \log_{12}(q) = \log_{16}(p+q)$


 What is the value of $\frac{q}{p}$?


 $\textbf{(A)}\ \frac{4}{3}\qquad \textbf{(B)}\ \frac{1+\sqrt{3}}{2}\qquad \textbf{(C)}\ \frac{8}{5}\qquad \textbf{(D)}\ \frac{1+\sqrt{5}}{2}\qquad \textbf{(E)}\ \frac{16}{9}$

 \textbf{Answer: }D



\textbf{Solution 1}

We can rewrite the equation as $\frac{\log{p}}{\log{9}} = \frac{\log{q}}{\log{12}} = \frac{\log{(p + q)}}{\log{16}}$. Then, the system can be split into 3 pairs: $\frac{\log{p}}{\log{9}} = \frac{\log{q}}{\log{12}}$, $\frac{\log{q}}{\log{12}} = \frac{\log{(p + q)}}{\log{16}}$, and $\frac{\log{p}}{\log{9}} = \frac{\log{(p + q)}}{\log{16}}$. Cross-multiplying in the first two, we obtain: 
 \[(\log{12})(\log{p}) = (2\log{3})(\log{q})\] and 
 \[(\log{12})(\log{(p + q)}) = (2\log{4})(\log{q})\]Adding these equations results in: 
 \[(\log{12})(\log{p(p+q)}) = (2\log{12})(\log{q})\] which simplifies to 
 \[p(p + q) = q^2\] Dividing by $pq$ on both sides gives: $\frac{p+q}{q} = \frac{q}{p} = \frac{p}{q} + 1$. We set the desired value, $q/p$ to $x$ and substitute it into our equation: $\frac{1}{x} + 1 = x \implies x^2 - x - 1 = 0$ which is solved to get our answer: $\boxed{\text{(D) } \frac{1 + \sqrt{5}}{2}}$. -lucasxia01

 



\textbf{Solution 2}

For some number t:

 $p = 9^{t}$
$q = 12^{t}$
$p + q = 16^{t}$

 Next we can divide $p + q$ by $p$ to obtain$\frac{p+q}{p} = 1 + \frac{q}{p}$

 Furthermore, we know that

 $\frac{p+q}{p} = (\frac{16}{9})^{t}$ and $\frac{q}{p} = (\frac{4}{3})^{t}$

 Substituting into the previous equation, we get $(\frac{16}{9})^{t} = 1 + (\frac{4}{3})^{t}$

 Let $x = (\frac{4}{3})^{t}$ and we can observe that $x = \frac{q}{p}$, then similarly to solution 1: $x^2 = 1 + x$, in which we get: $\boxed{\text{(D) } \frac{1 + \sqrt{5}}{2}}$ - ehmmaq

 


	\item (\textit{AHSME 1988}) In the figure, $AB \perp BC, BC \perp CD$, and $BC$ is tangent to the circle with center $O$ and diameter $AD$. 
In which one of the following cases is the area of $ABCD$ an integer?


 
                    \begin{figure}[H]
                        \centering
                        \includegraphics[width=0.5\textwidth]{Problem/AHSME_I_1988_27_Problem_Figure_1.png}
                    \end{figure}
                    
$\textbf{(A)}\ AB=3, CD=1\qquad \textbf{(B)}\ AB=5, CD=2\qquad \textbf{(C)}\ AB=7, CD=3\qquad\\ \textbf{(D)}\ AB=9, CD=4\qquad \textbf{(E)}\ AB=11, CD=5$

 \textbf{Answer: }D



\textbf{Solution}

Let $E$ and $F$ be the intersections of lines $AB$ and $BC$ with the circle. One can prove that $BCDE$ is a rectangle, so $BE=CD$.

 In order for the area of trapezoid $ABCD$ to be an integer, the expression $\frac{(AB+CD)BC}2=(AB+CD)BF$ must be an integer, so $BF$ must be rational.

 By Power of a Point, $AB\cdot BE=BF^2\implies AB\cdot CD=BF^2$, so $AB\cdot CD$ must be a perfect square. Among the choices, the only one where $AB\cdot CD$ is a perfect square is $\textbf{(D)}\ AB=9, CD=4$



	\item (\textit{AHSME 1988}) An unfair coin has probability $p$ of coming up heads on a single toss. 
Let $w$ be the probability that, in $5$ independent toss of this coin, 
heads come up exactly $3$ times. If $w = 144 / 625$, then 


 $\textbf{(A)}\ p\text{ must be }\frac{2}{5}\qquad \textbf{(B)}\ p\text{ must be }\frac{3}{5}\qquad\\ \textbf{(C)}\ p\text{ must be greater than }\frac{3}{5}\qquad \textbf{(D)}\ p\text{ is not uniquely determined}\qquad\\ \textbf{(E)}\ \text{there is no value of } p \text{ for which }w =\frac{144}{625}$

 \textbf{Answer: }D



\textbf{Solution}

We have $w = {5\choose3}p^{3}(1-p)^{2} = 10p^{3}(1-p)^{2}$, so we need to solve $10p^{3}(1-p)^{2} = \frac{144}{625} \implies p^{3}(1-p)^{2} = \frac{72}{3125}$. Now observe that when $p=0$, the left-hand side evaluates to $0$, which is less than $\frac{72}{3125}$; when $p=\frac{1}{2}$, it evaluates to $\frac{1}{32}$, which is more than $\frac{72}{3125}$, and when $p=1$, it evaluates to $0$ again. Thus, since $p^{3}(1-p)^{2}$ is continuous, the Intermediate Value Theorem tells us there is a solution between $0$ and $\frac{1}{2}$ and another solution between $\frac{1}{2}$ and $1$, meaning that there is not a unique value of $p$, so the answer is $\boxed{\text{D}}$.

 


 


	\item (\textit{AHSME 1988}) You plot weight $(y)$ against height $(x)$ for three of your friends and obtain the points 
$(x_{1},y_{1}), (x_{2},y_{2}), (x_{3},y_{3})$. If $x_{1} < x_{2} < x_{3}$ and $x_{3} - x_{2} = x_{2} - x_{1}$, 
which of the following is necessarily the slope of the line which best fits the data? 
"Best fits" means that the sum of the squares of the vertical distances from the data points to the line is smaller than for any other line.


 $\textbf{(A)}\ \frac{y_{3}-y_{1}}{x_{3}-x_{1}}\qquad \textbf{(B)}\ \frac{(y_{2}-y_{1})-(y_{3}-y_{2})}{x_{3}-x_{1}}\qquad\\ \textbf{(C)}\ \frac{2y_{3}-y_{1}-y_{2}}{2x_{3}-x_{1}-x_{2}}\qquad \textbf{(D)}\ \frac{y_{2}-y_{1}}{x_{2}-x_{1}}+\frac{y_{3}-y_{2}}{x_{3}-x_{2}}\qquad\\ \textbf{(E)}\ \text{none of these}$

 \textbf{Answer: }A



\textbf{Solution 1}

Let $y = mx + b$ a prediction line, and let $d = x_3 - x_2 = x_2 - x_1$.  Then $x_1 = x_2 - d$ and $x_3 = x_2 + d$.  Let $z_1 = mx_1 + b$, $z_2 = mx_2 + b$ and $z_3 = mx_3 + b$ be the predicted values of $y_1$, $y_2$, and $y_3$, respectively.  Then $z_1 = m(x_2 - d) + b = z_2 - md$ and $z_3 = m(x_2 + d) + b = z_2 + md$.

 Consider all lines that pass through $(x_2, z_2)$.  Let $r_1 = z_1 - y_1$, $r_2 = z_2 - y_2$, and $r_3 = z_3 - y_3$ be the three residuals.  Then $r_1 + r_3 = z_1 + z_3 - y_1 - y_3 = z_2 + md + z_2 - md - y_1 - y_3 = 2z_2 - y_1 - y_3$.  So $r_1 + r_3$ is the same for all these lines, as is $r_2$.  Let $s = r_1 + r_3$.  

 Since the best line passing through $(x_2, z_2)$ must minimize $r_1^2 + r_2^2 + r_3^2$, it must minimize $r_1^2 + r_3^2 = s^2 - 2r_1(s - r_1)$, and therefore maximize $r_1(s - r_1) = -r_1^2 + sr_1$.  Since this quadratic has a negative leading coefficient, it is maximized at $r_1 = -\frac{s}{2(-1)} = \frac{s}{2}$.  In this case, $r_1 = r_3 = \frac{s}{2}$, so $z_1 = y_1 + \frac{s}{2}$ and $z_3 = y_3 + \frac{s}{2}$.  Also, $z_3 - z_1 = y_3 - y_1$, and the slope of the line is $\frac{z_3 - z_1}{x_3 - x_1} = \frac{y_3 - y_1}{x_3 - x_1}$.

 Therefore, the line of best fit must have slope $\boxed{(\textbf{A}) \frac{y_3 - y_1}{x_3 - x_1}}$, as any line with a different slope could be improved by replacing it with the line with the same predicted value of $y_2$ and slope $\frac{y_3 - y_1}{x_3 - x_1}$.  

 -j314andrews

 



\textbf{Solution 2}

Apply one of the standard formulae for the gradient of the line of best fit, e.g. $\frac{\frac{\sum {x_i y_i}}{n} - \bar{x} \bar{y}}{\frac{\sum {x_{i}^2}}{n} - \bar{x}^2}$, and substitute in the given condition $x_3 - x_2 = x_2 - x_1$. The answer is $\boxed{\text{A}}$.

 



\textbf{Solution 3 (Linear Algebra)}

Let $A$ be a matrix and $x$ and $y$ be vectors. Recall that the least-squares error $||Ax - y||^2_2$ is minimized when $x = (A^T A)^{-1} A^T y$.

 In this case, $A = \begin{bmatrix}x_1 && 1 \\ x_2 && 1 \\x_3 && 1\end{bmatrix}$, $y = \begin{bmatrix}y_1\\y_2\\y_3\end{bmatrix}$, and $x = \begin{bmatrix}m \\ b\end{bmatrix}$, where $m$ and $b$ are the slope and $y$-intercept of the line of best fit, respectively.  Let $d = x_3 - x_2 = x_2 - x_1$, so $x_3 = x_2 + d$ and $x_1 = x_2 - d$.  Plugging these into $x = (A^T A)^{-1} A^T y$ yields $m = \frac{y_3 - y_1}{2d} = \boxed{(\textbf{A}) \frac{y_3 - y_1}{x_3 - x_1}}$.

 -j314andrews

 


	\item (\textit{AHSME 1988}) Let $f(x) = 4x - x^{2}$. Give $x_{0}$, consider the sequence defined by $x_{n} = f(x_{n-1})$ for all $n \ge 1$. 
For how many real numbers $x_{0}$ will the sequence $x_{0}, x_{1}, x_{2}, \ldots$ take on only a finite number of different values?


 $\textbf{(A)}\ \text{0}\qquad \textbf{(B)}\ \text{1 or 2}\qquad \textbf{(C)}\ \text{3, 4, 5 or 6}\qquad\\ \textbf{(D)}\ \text{more than 6 but finitely many}\qquad\\ \textbf{(E) }\infty$

 \textbf{Answer: }E



\textbf{Solution}

Note that $x_{0} = 0$ gives the constant sequence $0, 0, ...$, since $f(0) = 4 \cdot 0 - 0^2 = 0$. Because $f(4)=0, x_{0} = 4$ gives the sequence $4, 0, 0, ...$ with two different values. Similarly, $f(2) = 4$, so $x_{0} = 2$ gives the sequence $2, 4, 0, 0, ...$ with three values. In general, if $x_{0} = a_{n}$ gives the sequence $a_{n}, a_{n-1}, ... , a_{2}, a_{1}, ...$ with $n$ different values, and $f(a_{n+1}) = a_{n}$, then $x_{0} = a_{n+1}$ gives a sequence with $n+1$ different values. (It is easy to see that we could not have $a_{n+1} = a_{i}$ for some $i < n + 1$.) Thus, it follows by induction that there is a sequence with $n$ distinct values for every positive integer $n$, as long as we can verify that there is always a real number $a_{n+1}$ such that $f(a_{n+1}) = a_{n}$. This makes the answer $\boxed{\text{E}}$. The verification alluded to above, which completes the proof, follows from the quadratic formula: the solutions to $f(a_{n+1}) = 4a_{n+1} - a_{n+1}^{2} = a_{n}$ are $a_{n+1} = 2 \pm \sqrt{4 - a_{n}}$. Hence if $0 \leq a_{n} \leq 4$, then $a_{n+1}$ is real, since the part under the square root is non-negative, and in fact $0 \leq a_{n+1} \leq 4$, since $4-a_{n}$ will be between $0$ and $4$, so the square root will be between $0$ and $2$, and $2 \pm$ something between $0$ and $2$ gives something between $0$ and $4$. Finally, since $a_{1} = 0 \leq 4$, it follows by induction that all terms satisfy $0 \leq a_{n} \leq 4$; in particular, they are all real.

 


\end{enumerate}